% Единственный источник содержания темы. UTF-8.
% Метки label и ключи в [...] сохраняют опубликованные ссылки.

\section{Теория}

\subsection{Случайные события}

\begin{definition}
\textbf{Элементарное событие}  —  результат случайного опыта
\end{definition}

\textbf{!}Результатом случайного опыта всегда является только \textbf{одно} элементарное событие

\begin{definition}
\textbf{Множество элементарных событий}  —  множество, состоящее из всех элементарных событий
\end{definition}

\begin{designation}
\textbf{Множество элементарных событий}  —  $\Omega$
\end{designation}

\begin{definition}
\textbf{Случайное событие}  —  некоторое подмножество \emph{множества элементарных событий}
\end{definition}

\textbf{Пустое событие} $\emptyset$  —  случайное событие, не содержащее ни одного элементарного события.

\begin{definition}
\textbf{Невозможное} событие  —  пустое событие
\end{definition}

\begin{definition}
\textbf{Достоверное событие}  —  случайное событие, состоящее из всех элементарных событий (т.е. всё множество элементарных событий)
\end{definition}

\subsection{Противоположные события}

Случайные события будем обозначать большими латинскими буквами.\\
 Пусть нам дано некоторое случайное событие $A$.\\

\begin{definition}
\textbf{Противоположным событием} к событию $A$ будем называть дополнение $A$ к $\Omega$. То есть множество всех элементарных событий не лежащих в $A$
\end{definition}

\begin{designation}
\textbf{Противоположное событие} к событию $A$  —  $\overline{A}$
\end{designation}

$\overline{A} = \Omega \backslash A$. Обратный слэш $\backslash$ между множествами означает разность множеств.\\
 $A\cup\overline{A} = \Omega$\\
 Если $\overline{A}$  —  событие, противоположное $A$, то $A$  —  событие, противоположное $\overline{A}$. Такие события называют \textbf{взаимнопротивоположными}.

\figure{opposite-events.png}{Событие A и противоположное событие в пространстве исходов}{1860}{1860}

\subsection{Вероятность случайного события}

Неформально: \textbf{Вероятность}  —  число, оценивающее возможность наступления данного события. Каждому событию мы будем сопоставлять некоторое вещественно число из отрезка $[0,1]$. Причём достоверному событию всегда соответствует $1$, а невозможному  —  $0$. Сопоставляемое число и будет называться \textbf{вероятностью случайного события}\\

\begin{designation}
\textbf{Вероятность}  —  $\mathbb{P}$ или $\mathbf{P}$
\end{designation}

Для тех кому интересно, ниже более формальное определение вероятности (учить аксиомы не нужно):\\

\begin{definition}
\textbf{Вероятность (вероятностная мера)}  —  функция из множества событий в отрезок $[0,1]$, удовлетворяющая трём аксиомам
\end{definition}

\begin{axioms}[random-events/theorem/1]\label{theorem-1}
\textbf{Аксиомы вероятности}:\\

\begin{enumerate}
\item \emph{Неотрицательность.} Вероятность любого события неотрицательна $\mathbb{P}(A) \geq 0$
\item \emph{Конечность.} Вероятность достоверного события равна $1$. $\mathbb{P}(\Omega) = 1$
\item \emph{Аддитивность.} Для двух \textbf{несовместных} событий $A, B$ вероятность объединения равна сумме вероятностей каждого из событий. $A\cap B = \emptyset, \mathbb{P}(A\cup B) = \mathbb{P}(A) + \mathbb{P}(B)$
\end{enumerate}
\end{axioms}

\begin{explanation}[random-events/theorem/1]\label{solution-random-events-theorem-1}
Аксиомы — исходные свойства, с помощью которых определяется вероятность. В этой модели их принимают без доказательства. Теоремы и правила вычисления затем выводят из них; например, сумма вероятностей элементарных событий равна 1 вследствие аддитивности и равенства $\mathbb P(\Omega)=1$.

В этом блоке перечислены исходные предположения вероятностной модели.
\end{explanation}

P.s. Про определение \textbf{несовместных событий} и сложение вероятностей будет дальше.

\subsection{Вероятность элементарного события и равновозможные события}

Почти всегда в школе мы будем обсуждать случаи, когда множество элементарных событий будет конечным.\\

\begin{theorem}[random-events/theorem/2]\label{theorem-2}
Сумма вероятностей всех элементарных событий равна $1$.
\end{theorem}

\begin{proof}[random-events/theorem/2]\label{solution-random-events-theorem-2}
В рассматриваемом конечном пространстве исходов обозначим все элементарные события через $\omega_1,\ldots,\omega_N$. Они попарно несовместны и в объединении дают достоверное событие $\Omega$.

Последовательно применяя аддитивность к объединению несовместных событий, получаем

\[\mathbb P(\omega_1)+\cdots+\mathbb P(\omega_N)=\mathbb P(\Omega)=1.\]
\end{proof}

\begin{definition}
Если вероятности всех элементарных событий равны, то говорят, что это случайный опыт с \textbf{равновозможными элементарными событиями}
\end{definition}

\begin{theorem}[random-events/theorem/3]\label{theorem-3}
При опыте с $N$ \emph{равновозможными элементарными событиями} вероятность каждого элементарного события равна $\frac{1}{N}$
\end{theorem}

\begin{proof}[random-events/theorem/3]\label{solution-random-events-theorem-3}
По равновозможности вероятность каждого из $N$ элементарных событий имеет одно и то же значение $p$. По предыдущей теореме сумма этих вероятностей равна 1, то есть $Np=1$. Так как $N\ge1$, получаем $p=\frac1N$.
\end{proof}

\section{Задачи}

\begin{task}[random-events/problems/1]\label{task-problems-1}
\textbf{Задача 1.} Пусть множество элементарных событий состоит из $N$ элементарных событий. Найдите вероятность каждого элементарного события.
\end{task}

\begin{solution}[random-events/problems/1]\label{solution-random-events-problems-1}
Обозначим общую вероятность одного элементарного события через $p$. Все $N$ элементарных событий несовместны, а их объединение — всё пространство исходов. Поэтому $Np=1$, откуда $p=\frac1N$.

\textbf{Ответ:} вероятность каждого элементарного события равна $\frac1N$.
\end{solution}

\begin{task}[random-events/problems/2]\label{task-problems-2}
\textbf{Задача 2. Ошибка Д'Аламбера} Симметричную монету бросают дважды. Перечислите все элементарные события, найдите вероятность каждого из них. Теперь две одинаковые симметричные монеты бросают одновременно. Перечислите все элементарные события и найдите вероятность каждого из них.
\end{task}

\begin{solution}[random-events/problems/2]\label{solution-random-events-problems-2}
Обозначим орёл буквой О, решку — Р. При последовательных бросках различаем первый и второй результаты:

\[\Omega=\{\text{ОО},\text{ОР},\text{РО},\text{РР}\}.\]

Все четыре исхода равновозможны, вероятность каждого равна $\frac14$.

При одновременном броске можно мысленно пометить монеты номерами 1 и 2: получатся те же четыре равновозможных исхода. Если номера монет не учитывать, наблюдаются три результата: два орла, один орёл и одна решка, две решки. Средний результат объединяет два помеченных исхода — ОР и РО, поэтому вероятности этих трёх результатов равны соответственно $\frac14$, $\frac12$, $\frac14$.

Ошибка Д'Аламбера состоит в том, чтобы считать эти три наблюдаемых результата равновозможными и приписывать каждому вероятность $\frac13$.
\end{solution}

\begin{task}[random-events/problems/3]\label{task-problems-3}
\textbf{Задача 3. №259} Игральную кость подбрасывают несколько раз. Равновозможны ли элементарные события этого опыта. Найдите вероятность каждого элементарного события при\\
 а) двух бросках\\
 б) трёх бросках\\
 в) $n$ бросках
\end{task}

\begin{solution}[random-events/problems/3]\label{solution-random-events-problems-3}
В стандартной модели симметричной игральной кости и независимых бросков элементарный исход — упорядоченная последовательность выпавших чисел. Для каждого следующего результата есть 6 вариантов.

а) При двух бросках получаем $6\cdot6=36$ равновозможных исходов; вероятность каждого — $\frac1{36}$.

б) При трёх бросках исходов $6^3=216$; вероятность каждого — $\frac1{216}$.

в) При $n$ бросках исходов $6^n$; вероятность каждого — $\frac1{6^n}$.
\end{solution}

\section{Домашнее задание}

Стр. 142–143.

\begin{task}[random-events/homework/1]\label{task-homework-1}
\textbf{1. №249 (г)} Случайный опыт может закончиться одним из трёх элементарных событий: $a$, $b$ или $c$. Чему равна вероятность элементарного события $c$, если $\mathbb{P}(a)=p$, $\mathbb{P}(b)=0,8-p$? Какие значения может принимать $p$?
\end{task}

\begin{solution}[random-events/homework/1]\label{solution-random-events-homework-1}
Вероятности трёх элементарных событий в сумме дают 1:

\[\mathbb P(c)=1-\mathbb P(a)-\mathbb P(b)=1-p-(0,8-p)=0,2.\]

Неотрицательность двух остальных вероятностей даёт $p\ge0$ и $0,8-p\ge0$. Поэтому $0\le p\le0,8$. При любом таком $p$ все три вероятности лежат в $[0,1]$ и в сумме равны 1.

\textbf{Ответ:} $\mathbb P(c)=0,2$, $p\in[0;0,8]$.
\end{solution}

\begin{task}[random-events/homework/2]\label{task-homework-2}
\textbf{2. №250} Игральная кость несимметрична. В таблице 49 показаны вероятности выпадения на этой кости 1, 2, 3, 5 или 6 очков. Найдите вероятность выпадения 4 очков.
\end{task}

\begin{tabledata}{Таблица 49. Вероятности выпадения граней}
Число очков & 1 & 2 & 3 & 4 & 5 & 6 \\
Вероятность & $\frac{1}{4}$ & $\frac{1}{12}$ & $\frac{1}{6}$ & ? & $\frac{1}{12}$ & $\frac{1}{4}$ \\
\end{tabledata}

\begin{solution}[random-events/homework/2]\label{solution-random-events-homework-2}
Сумма вероятностей выпадения всех шести граней равна 1. Обозначим неизвестную вероятность через $x$:

\[x=1-\left(\frac14+\frac1{12}+\frac16+\frac1{12}+\frac14\right)=1-\frac{10}{12}=\frac16.\]

\textbf{Ответ:} вероятность выпадения четырёх очков равна $\frac16$. Это не означает, что вся кость симметрична: вероятности остальных граней различаются.
\end{solution}

\begin{task}[random-events/homework/3]\label{task-homework-3}
\textbf{3. №257} Три первоклассника по очереди выбирают воздушные шарики. Каждый из них выбирает шарик одного из двух цветов: зелёного (З) или синего (С). Выпишите элементарные события этого эксперимента. Считая, что все они равновозможны, найдите вероятность каждого из них.
\end{task}

\begin{solution}[random-events/homework/3]\label{solution-random-events-homework-3}
Записываем цвета в порядке выбора детьми. Например, ЗСЗ означает: первый ребёнок выбрал зелёный шарик, второй — синий, третий — зелёный.

Все элементарные события: ЗЗЗ, ЗЗС, ЗСЗ, ЗСС, СЗЗ, СЗС, ССЗ, ССС. Их $2^3=8$. По условию они равновозможны, поэтому вероятность каждого равна $\frac18$.
\end{solution}

\begin{task}[random-events/homework/4]\label{task-homework-4}
\textbf{4. №258} Три первоклассника по очереди выбирают фломастеры. Каждый из них выбирает фломастер одного из трёх цветов: зелёного (З), синего (С) или красного (К). Сколько у этого опыта элементарных событий? Считая, что все элементарные события равновозможны, найдите вероятность каждого из них.
\end{task}

\begin{solution}[random-events/homework/4]\label{solution-random-events-homework-4}
Элементарный исход — упорядоченная тройка цветов. Каждый из трёх детей выбирает один из трёх цветов, поэтому всего $3\cdot3\cdot3=27$ исходов. По условию они равновозможны.

\textbf{Ответ:} 27 элементарных событий; вероятность каждого — $\frac1{27}$.
\end{solution}

